\documentclass[../main.tex]{subfiles}
\graphicspath{{\currfiledir}}
\begin{document}
\chapter{Terminology}

\section{Grapheme}
\textcquote[204]{crystal2010}{Graphemes are the smallest units in a writing system capable of 
causing a contrast in meaning.}
In linguistics, graphemes are often placed into angle brackets, e.g. \grapheme{a} or \grapheme{b}.
Sometimes \emph{graphemes} are called \emph{signs}.
The term \emph{sign} and \emph{grapheme} are considered to be equivalent and 
interchangeable throughout this work.

\section{Allograph}
\textcquote[204]{crystal2010}{Graphemes are abstract units which may adopt a variety of forms 
\elide Each of these possible forms is known as \emph{graphs}\elide
There is a vast amount of physical variation in the shapes of graphs that does not affect the 
underlying identity of the grapheme\elide 
When graphs are analyzed as variants of a grapheme, they are known as \emph{allographs}.}
The Maya script uses a lot of allographs.
For example, the syllable \syllable{u} can be written in many ways all having the same meaning.
See \ref{fig:terminology-grapheme-u-allographs} for a selection of allographs.
\begin{center}
    \includegraphics[width=\textwidth,keepaspectratio]{img/grapheme-u-allographs}
    \captionof{figure}{Some allographs of the grapheme \grapheme{u}}
    \label{fig:terminology-grapheme-u-allographs}
\end{center}

\section{Glyph and glyph block}
The term \emph{Hieroglyph} or \emph{glyph} is not an exact term.
In epigraphic literature, a \emph{glyph} is considered to be a group of one or more graphemes.
They can represent a syllable, single words or even a whole phrase 
(\cite[1]{bricker1986}, (\cite[23]{macrilooper2003})).
\emph{Glyphs} are usually arranged in regular rectangular shapes called a \emph{glyph block}.
\textcquote[23]{macrilooper2003}{The rectangular shape of \emph{glyph blocks} results from 
the arrangement of texts into rows and columns}.

\section{Analysis of an unknown writing system}
In distinguishing distinctive signs and their allographs, prior to linguistic decipherment, 
researchers followed the assumption that graphemes of a script are considered the same if 
they resemble each other in more features than either resembles any other.
The convenient way to do this is the segmentation of texts into characters and the subsequent 
verification by means of distribution analyses, which establishes the distribution and 
connection among one another and confirms identification. 
If there is no linguistic content control (character interpretation), 
however, there is the uncertainty that graphemes and 
their allographs are not recognized in writing systems rich in variants and erroneously 
inventoried as distinctive characters with their own number key.


Yuri Knorozov states that, \textcquote[34]{knorozov1967}{Two [signs] are identical when 
they are both composed of the same graphic elements\elide, whose drawing and disposition is
sufficiently similar to allow them to be identified.}



Identifying graphemes in an unknown writing system can be challenging.




Eric Thompson (\cite[12\psq]{thompson1962catalog}) also recognized the method as a source of 
potential false conclusions.


Tikal Temple IV, Lintel 2 A7


\end{document}
